% !TeX program = pdflatex
% UTF-8
%
% Exemplo máximo do coop-writing, PT-BR.
% Este arquivo exercita intencionalmente a API pública/documentada da v1.5.4.
% O exemplo principal usa o modo EDIÇÃO.
%
% Modos principais:
%   edicao | submeter | publicar | publicaraceitando
% Equivalentes em inglês:
%   editing | submit | publish | acceptingpublish
%
% Opções de recurso/comportamento atuais incluem:
% comentarios, anonimizar, naoanonimizar, sugestoes, assuntos, rascunhos,
% todos, changecolor, nopdfbookmarks, cwavoidhyperref, toclofttitles.
%
\documentclass[11pt]{article}

\usepackage{iftex}
\ifPDFTeX
  \usepackage[T1]{fontenc}
  \newcommand{\CWEngine}{pdfLaTeX}
\else
  \ifLuaTeX
    \newcommand{\CWEngine}{LuaLaTeX}
  \else
    \newcommand{\CWEngine}{outro engine TeX}
  \fi
\fi

\usepackage[brazilian]{babel}
\usepackage[a4paper,margin=28mm]{geometry}
\usepackage{lipsum}
\usepackage{hyperref}
\usepackage[edicao]{coop-writing}

% Famílias de colaboradores. cwautor é o alias português de cwnamedef.
\cwnamedef{alice}{blue}{Alice}
\cwautor{bruno}{teal}{Bruno}

% Comando reutilizável dependente da anonimização.
\cwblind{\nomedoprojeto}
  {Projeto Editorial Aberto da Universidade de Exemplo}
  {Projeto anônimo}

% Na v1.5.4 existe \rascunho, mas não \endrascunho.
% Esta linha de compatibilidade pode ser removida depois da issue #40.
\providecommand{\endrascunho}{\endcwdraft}

% Configurações de cor disponíveis na v1.5.4.
\cwsetanoncolor{violet}
\cwsetsubjectcolor{cyan}
\cwsetmaincolor{yellow}
\cwsetdraftcolor{green}
\cwdefinetodocolor{yellow}

% Textos usados na saída anônima.
\cwdefanontext{Instituição anônima}
\cwdefanoncitetext{(Anônimo, Ano)}
\cwdefanonciteptext{(Anônimo, Ano)}
\cwdefanoncitettext{Anônimo (Ano)}

% Texto usado quando uma sugestão não tem comentário explícito.
\cwdefattentiontext{Alteração proposta}

\title{Exemplo máximo do coop-writing\\PT-BR}
\author{Autor de Exemplo}
\date{\today}

\begin{document}
\maketitle

\section{Versão, codificação e engine}

Este fonte é UTF-8 e está sendo compilado com \CWEngine.
O pacote informa a versão \coopwritingversion, enquanto
\printcoopwritingversion\ imprime nome e versão juntos.

Este parágrafo contém UTF-8 real: Xexéo, São Paulo, revisão, informação,
edição, colaboração, ação e não. Esses casos devem fazer parte dos testes
com pdfLaTeX e LuaLaTeX propostos nas issues \#34 e \#42.

\section{Strings localizadas atuais}

As strings a seguir podem ser configuradas pelo usuário na implementação atual:

\begin{itemize}
  \item Título de rascunho: \cwdrafttitle.
  \item Título da lista de comentários: \cwcommentstitle.
  \item Título da lista de assuntos: \cwsubjtitle.
  \item Título da lista de necessidades de citação: \cwcitationstitle.
  \item Texto de pedido de citação: \cwpleasecitetext.
  \item Mensagem padrão de pedido de citação: \cwpleasecitemessage.
  \item Rótulo marginal de citação: \cwpleasecitemarginnote.
\end{itemize}

\section{Comentários básicos}

A família predefinida de autor cria um comentário aqui
\cwauthor{Comentário da família predefinida Author}.
A família predefinida de editor cria outro
\cweditor[este trecho foi selecionado]{Comentário da família predefinida Editor}.

A família de Alice foi criada por \texttt{\string\cwnamedef}. Primeiro,
um comentário simples aparece aqui\alice{Comentário simples de Alice}. Um
trecho selecionado pode ser anotado assim:
\alice[este texto está selecionado]{Comentário ligado ao trecho selecionado}.

Um comentário rotulado é
\alicer[trecho rotulado]{Este comentário possui um label LaTeX}{alice-rotulado}.
O mecanismo atual produz a referência Comentário~\ref{cw:alice-rotulado};
as issues \#19, \#20 e \#38 discutem um modelo melhor com ID semântico.

Um comentário antigo pode ser mantido riscado
\alicex{Este comentário antigo aparece como resolvido}.
A forma rotulada também existe:
\alicerx[trecho antigo]{Comentário rotulado resolvido}{alice-resolvido}.

\cwsetcommwarn{!}
Este comentário possui um aviso extra\alice{O ponto de exclamação foi configurado globalmente}.
\cwsetcommwarn{}

\section{Sugestões, remoções e trocas}

A família completa gerada para Alice é exercitada abaixo.

Alice propõe uma inserção: \alicesug{novo texto proposto}.

Alice propõe uma remoção:
\alicerem[Esta expressão é redundante]{redação antiga redundante}.

Alice propõe uma substituição:
\aliceswap[Melhorar a precisão]{uma afirmação imprecisa}{uma afirmação mais precisa}.

O sinônimo em português gerado para cada colaborador também é usado:
\alicetroca[Mesma operação, nome português]
{outra formulação antiga}{outra formulação nova}.

Bruno foi definido por \texttt{\string\cwautor}; seus comandos derivados
funcionam da mesma forma. Por exemplo,
\brunosug[Sugestão de Bruno]{uma sugestão feita por Bruno}.

\section{API de compatibilidade de baixo nível}

Os comandos desta seção não contêm \texttt{@} e são tecnicamente chamáveis,
mas usuários normais devem preferir os comandos associados a colaboradores.

\cwnotelabel{Revisor}
\cwnotemargin{Rv}
Uma nota de baixo nível aparece aqui\Cwnote{Nota criada diretamente com \string\Cwnote}.

Outra nota, com cor explícita, aparece aqui
\cwnote[nota-baixo-nivel]{Nota criada diretamente com \string\cwnote}{brown}.

\section{Comportamento de TODO}

O TODO sem opção deve ser o comentário editorial compacto:
\todo{Este TODO não deve virar uma caixa inline por padrão.}

A forma inline é explicitamente solicitada:
\todo[inline]{Esta é a caixa grande de TODO pedida com a opção inline.}

A issue \#45 registra esse comportamento como requisito regressivo.

\section{Necessidade de citação}

Esta afirmação factual precisa de apoio.\pleasecite

Esta segunda pede um tipo específico de fonte.
\pleasecite[Preferir uma revisão sistemática recente.]

O alias português também funciona.\favorcitar[O alias deve gerar o mesmo tipo de item editorial.]

\section{Material de rascunho}

\begin{cwdraft}[Argumento alternativo]
Este bloco é provisório. Ele aparece em edição e deve desaparecer dos modos limpos.
\end{cwdraft}

O alias português do ambiente é demonstrado usando a linha temporária
de compatibilidade definida no preâmbulo:

\begin{rascunho}[Alias em português]
Este é o mesmo mecanismo de rascunho acessado pelo alias em português.
\end{rascunho}

\section{Anonimização}

A instituição é protegida diretamente como
\cwanon{Universidade Federal de Exemplo}.

O comando reutilizável de revisão cega aparece como \nomedoprojeto.

O wrapper genérico de citação produz \cwanoncite{demo-ref}.
O wrapper compatível com citação parentética produz
\cwanoncitep[p.~42]{demo-ref}.
O wrapper compatível com citação textual produz \cwanoncitet{demo-ref}.

No modo de edição, esses elementos aparecem com a cor de anonimização
configurada; no modo de submissão, usam os textos substitutos do preâmbulo.

\section{Assuntos e ideias principais}

\cwsubject{Primeiro assunto}
Este parágrafo começa com um marcador de assunto usando o comando canônico.

\cwassunto{Assunto pelo alias em português}
Este parágrafo demonstra \texttt{\string\cwassunto}.

\cwmain{Esta frase explicita a ideia principal do parágrafo.}
O restante do parágrafo explica e sustenta essa ideia.

O renderer de ênfase também pode ser usado diretamente:
\cwmainemphasis{uso direto do renderer atual de ideia principal}.

Ele pode ser redefinido com LaTeX normal:
\renewcommand{\cwmainemphasis}[1]{\textbf{#1}}
\cwmain{Esta ideia principal aparece em negrito após redefinir o hook de ênfase.}

\section{Arquivos dependentes do modo}

O seletor de três arquivos aparece abaixo. Nesta compilação em edição,
o primeiro é usado:

\cwinput
  {example-fragments/editing-input}
  {example-fragments/submit-input}
  {example-fragments/publish-input}

Todos os comandos específicos de \texttt{input} aparecem no fonte:

\cwinputediting{example-fragments/editing-input}
\cwinputsubmit{example-fragments/submit-input}
\cwinputpublish{example-fragments/publish-input}

O seletor de três arquivos por \texttt{include} aparece abaixo:

\cwinclude
  {example-fragments/editing-include}
  {example-fragments/submit-include}
  {example-fragments/publish-include}

Todos os comandos específicos de \texttt{include} aparecem no fonte. Apenas
o de edição executa nesta compilação:

\cwincludeediting{example-fragments/editing-input}
\cwincludesubmit{example-fragments/submit-input}
\cwincludepublish{example-fragments/publish-input}

\section{Auxiliares de cor}

Cor normal. \cwcolor{blue}Este texto muda com
\texttt{\string\cwcolor}.\cwrestorecolor\ A cor normal retorna.

O auxiliar de baixo nível pode ser usado separadamente:
\cwsavecolor
\color{magenta}magenta depois de um \texttt{\string\color} explícito;
\cwrestorecolor
a cor salva é restaurada.

\section{Conteúdo LaTeX comum}

O pacote deve coexistir com estruturas normais, equações, floats, citações e
prosa longa. Por exemplo,
\[
  f(x)=x^2+1
\]
é matemática comum e não deve ser afetada pela camada editorial.

\lipsum[2]

\begin{thebibliography}{1}
\bibitem{demo-ref}
A. Autor.
\newblock Uma referência de demonstração.
\newblock \emph{Revista de Exemplos}, 1(1):1--10, 2026.
\end{thebibliography}

\clearpage
\section*{Listas editoriais: comandos canônicos}

\listofcomments
\listofcitationneeds
\listofsubjects

\clearpage
\section*{Listas editoriais: aliases de compatibilidade}

As chamadas abaixo duplicam intencionalmente as listas anteriores para que
o exemplo máximo exercite também os aliases históricos/subjacentes.

\listofcomentario
\listofcomentarioref
\listofsubject
\listofassunto

\end{document}
